\newcommand{\premessa}{
    \chapter*{Premessa}
    \addcontentsline{toc}{chapter}{Premessa}
    
    Ho preparato l'esame di Interazione \& Multimedia senza seguire le lezioni in presenza, ma studiando in modo autonomo. Questo documento non intende in alcun modo sostituire le lezioni dei professori né i loro materiali di studio ufficiali.
    
    \vspace{0.5cm}
    
    Come recita il \textit{Capybara Compendium}:
    
    \begin{quote}
        \itshape
        "L'obbiettivo di questi riassunti è permettere un veloce ripasso degli argomenti trattati in uno specifico anno accademico, non vogliono in alcun modo sostituire le lezioni e i libri di testo consigliati. I contenuti non sono in alcun modo né approvati né revisionati dal prof. della relativa materia, sono appunti elaborati da uno studente, ergo qualità e correttezza non garantità nonostante il mio impegno."
    \end{quote}
    
    \vspace{0.5cm}
    
    \section*{Fonti utilizzate}
    
    Per la realizzazione di questa dispensa mi sono servito di:
    
    \subsection*{Parte Teorica}
    \begin{itemize}
        \item Slide di teoria del \textbf{Professore Filippo Stanco}
        \item Dispense di \textbf{MMorana094} (\href{https://github.com/MMorana094/Interazione-e-Multimedia}{https://github.com/MMorana094/Interazione-e-Multimedia})
        \item Dispense di \textbf{Rosario Terranova} (\href{https://github.com/rosarioterranova}{https://github.com/rosarioterranova})
        \item \textbf{Capybara Compendium} (\href{https://github.com/LightDestory/Capybara-Compendium}{https://github.com/LightDestory/Capybara-Compendium})
    \end{itemize}
    
    \subsection*{Parte Laboratorio (Processing)}
    \begin{itemize}
        \item Videolezioni del \textbf{Professore Lorenzo Di Silvestro}
        \item Documentazione ufficiale di \textbf{Processing} (\href{https://processing.org/reference}{https://processing.org/reference})
    \end{itemize}
}